\section*{Zusammenfassung}

Beweiszertifikate können das Vertrauen in automatisierte Deduktionswerkzeuge erhöhen, indem sie eine Begründung ihrer Schlussfolgerungen liefern welche von einem unabhängigen Prüfprogramm validiert werden kann.
Der Practical Algebraic Calculus (PAC) definiert solche Beweiszertifikate für die Verifikation digitaler Schaltungen, insbesondere von Multiplizierern, auf Grundlage algebraischer Methoden.
Ich migriere den formal verifizierten PAC Prüfer Pastèque von seinem ursprünglichen Standard-ML-Backend auf ein neueres LLVM-Backend und verwendete dafür das Isabelle-LLVM-Framework das die Implementierung und formale Verifikation imperativer Programme mit manueller Speicherverwaltung mit dem Theorembeweiser Isabelle ermöglicht.
Im Zuge dieser Migration erweiterte ich eine native Bibliothek für Mehrfachpräzisionsarithmetik in Isabelle-LLVM für die Verwendung in Pastèque und implementierte grundlegende Datentypen für Isabelle-LLVM, darunter Hashmaps und schachtelbare verkettete Listen.
Ich evaluierte die Performanz der LLVM-Migration im Vergleich zur ursprünglichen Standard-ML-Implementierung und zum in C++ geschriebenen PAC-Prüfer Pacheck. Dabei zeigt sich eine Leistungsverbesserung bei kleineren Instanzen und ein Leistungsrückgang bei großen 512-Bit-Multiplizierern, im Vergleich zum Standard-ML-Backend
